\addcontentsline{toc}{section}{Préface}
\section*{Préface}

C'est en me baladant sur Steam que j'ai découvert les \acrlong{vn}s, en tombant par hasard sur \textit{Higurashi When They Cry Hou - Ch.1 Onikakushi}. Ce qui m'a frappé, c'est la puissance narrative que ce genre peut atteindre avec des moyens apparemment simples: des images fixes, du texte, de la musique. Depuis, j'ai continué à explorer ce genre, avec des titres comme \textit{Marco et le dragon-galaxie} ou \textit{Find Love or Die Trying}, et mon envie de créer ma propre histoire n'a fait que grandir.

Lorsque je me suis penché sur les outils existants, j'ai rapidement déchanté. RenPy impose l'apprentissage d'un langage de script qui lui est propre. Twine demande une installation et une prise en main qui rebutent quiconque veut simplement écrire une histoire. Je ne voulais pas apprendre un énième outil, je voulais me concentrer sur le récit.

C'est cette frustration qui est à l'origine de \textit{NamelessStory}. L'idée initiale était simple, créer une \gls{librairie} \gls{reactjs} personnelle pour faire tourner mon propre \acrlong{vn}. Mais au fil du développement, le projet a évolué vers quelque chose de plus ambitieux, un moteur \gls{opensource}, accessible depuis un navigateur, pensé pour que n'importe qui disposant de notions de base en \acrshort{json} puisse créer et publier son histoire sans installer quoi que ce soit.

Ce mémoire documente cette évolution, d'une idée personnelle vers un outil conçu pour être partagé.